\chapter*{使用说明}
\addcontentsline{toc}{chapter}{使用说明}

本书是《PIE-Block 软件与应用》的简明版，面向首次使用 PIE-Block 的参赛队员与机械组成员。
完整版教材中的程序结构、编译产物与实现细节不在本书范围内；本书只保留从新建项目到实车
验证的必要步骤，以便在装车现场快速查阅。

全书先给出通用设置，再按项目类型组织。第~\ref{ch:concise-intro}~章介绍软件定位、
四种项目类型与通用流程；第~\ref{ch:concise-debug}~章用调试项目走通第一条
“配置—检查—编译—烧录”链路；第~\ref{ch:concise-remote}~章介绍遥控器、四轮底盘与
摩擦轮转速的通用设置；第~\ref{ch:concise-infantry}~章配置步兵；
第~\ref{ch:concise-engineer}~章配置工程；第~\ref{ch:concise-music}~章使用音乐项目。
附录汇总检查、烧录、常见问题与装车前检查单。

操作层面的三点约定：

\begin{itemize}
  \item 涉及实车的内容均在车体架空、速度限低、射界无人的前提下进行；
  \item 烧录前关闭主控板 SERVO、POWER 与机械拓展板 POWER、BOOSTER 四个供电开关；
  \item 界面截图取自当前版本，具体文案可能随版本略有调整，以软件中的实际显示为准。
\end{itemize}

\noindent PIE-Block 项目组
